\problem{CT01}{二阶概率的无偏估计}{充分统计量；Rao--Blackwell}
$X_1,\ldots,X_n$ 独立服从 $\operatorname{Bernoulli}(p)$，$n\ge2$、$0<p<1$。从无偏估计 $X_1X_2$ 出发，构造 $p^2$ 的方差更小或相等的无偏估计，求其精确方差，并证明它在所有无偏估计中方差最小。
\solution 令 $S=\sum_iX_i$。给定 $S=s$，成功位置在所有大小为 $s$ 的子集中等概率分布，故
\[
 U:=\E[X_1X_2\mid S]=\frac{S(S-1)}{n(n-1)}.
\]
条件分布不依赖 $p$，所以 $S$ 是充分统计量。全方差公式给出 $\Var(X_1X_2)=\Var(U)+\E\Var(X_1X_2\mid S)$，保证无偏性不变且方差不增加。
还可将 $U$ 写成所有 $\binom n2$ 个乘积 $X_iX_j$ 的平均。单项方差为 $p^2(1-p^2)$，共享一个指标的两项协方差为 $p^3(1-p)$，不相交两项独立。共享指标的无序项对有 $n\binom{n-1}{2}$ 个，因此
\[
 \Var(U)=\frac{2p^2(1-p^2)+4(n-2)p^3(1-p)}{n(n-1)}.
\]
为证明最小方差，先证完备性：若 $\E_p h(S)=0$ 对所有 $p\in(0,1)$ 成立，除以 $(1-p)^n$ 后得到关于 $p/(1-p)$ 的恒零多项式，所以每个系数 $\binom ns h(s)$ 都为零。
对任意另一无偏估计 $W$，其条件均值 $\E[W\mid S]$ 是不含未知参数的函数，且与 $U$ 的期望相同；完备性迫使二者相等。因此再次用全方差公式得 $\Var(W)\ge\Var(U)$。无有限方差的无偏估计也不可能改善该界。$n=2$ 时 $U=X_1X_2$，方差公式仍成立。
\insight 条件化去掉与参数无关的随机性；完备性进一步保证没有其他无偏的充分统计量函数可供改进。

\problem{CT02}{指数样本的最强检验}{似然比；精确功效}
$X_1,\ldots,X_n$ 独立服从率参数为 $\theta$ 的指数分布。检验 $H_0:\theta=\theta_0$ 对 $H_1:\theta=\theta_1>\theta_0$。在第一类错误概率至多 $\alpha\in(0,1)$ 的检验中，求最强检验及其功效。给出 $n=2,\theta_0=1,\theta_1=2,\alpha=1-2/e$ 时的明确临界值与功效。
\solution 令 $T=\sum_iX_i$。备择与原假设的联合密度比为
\[
 L=\left(\frac{\theta_1}{\theta_0}\right)^n
      e^{-(\theta_1-\theta_0)T},
\]
随 $T$ 严格递减，故候选检验为 $\phi_*=\ind_{\{T<c\}}$，其中 $c>0$ 唯一满足 $\Prob_{\theta_0}(T<c)=\alpha$。$T$ 连续，无须在边界随机化。
设相应似然比阈值为 $k$。任何其他检验 $0\le\phi\le1$ 都有 $(\phi_* -\phi)(L-k)\ge0$。若其大小不超过 $\alpha$，则
\[
 \E_{\theta_1}(\phi_* -\phi)
 \ge k\E_{\theta_0}(\phi_* -\phi)\ge0.
\]
所以候选检验的功效最大，这也给出了似然比原则的证明。
指数变量之和的密度为 $\theta^nt^{n-1}e^{-\theta t}/(n-1)!$；可由卷积归纳得到。逐次分部积分给出
\[
 \Prob_\theta(T<c)=1-e^{-\theta c}
       \sum_{j=0}^{n-1}\frac{(\theta c)^j}{j!}.
\]
用 $\theta_0$ 解临界值、再用 $\theta_1$ 代入即为功效。指定数值下 $c=1$，功效为 $1-3e^{-2}$。率参数越大，等待总时间倾向越短，因此拒绝域位于左尾。
\insight 先用似然比确定应选择哪一侧，再用原假设分布校准大小；检验方向不能只凭参数名称猜测。

\problem{CT03}{正态均值的收缩风险}{Stein 恒等式；平方损失}
观测 $X\sim N_d(\theta,I_d)$，$d\ge3$，未知 $\theta\in\R^d$。在估计族
\[
 \delta_a(X)=\left(1-\frac a{\|X\|^2}\right)X
\]
中求使平方损失风险最小的常数 $a$，并证明它对每个 $\theta$ 都优于直接使用 $X$。在 $X=0$ 处任意定义估计，不影响风险。若用同一估计预测独立的 $Y\sim N_d(\theta,I_d)$，其预测风险如何变化？
\solution 写 $\delta_a=X+g(X)$，其中 $g(x)=-ax/\|x\|^2$。对高斯密度分部积分可得
\[
 \E[(X_i-\theta_i)g_i(X)]=\E[\partial_i g_i(X)].
\]
本题的零点奇性不会破坏等式：先在 $\|x\|>\varepsilon$ 积分，球面边界项的量级为 $\varepsilon^{d-2}\to0$；导数的局部量级为 $\|x\|^{-2}$，在 $d>2$ 时可积，无穷远边界由高斯衰减消失。
直接求导得到
\[
 \|g(x)\|^2=\frac{a^2}{\|x\|^2},\qquad
 \sum_i\partial_i g_i(x)=-\frac{a(d-2)}{\|x\|^2}.
\]
展开风险并应用上述恒等式，得
\[
 \begin{aligned}
 R_\theta(a)&=\E_\theta\|\delta_a(X)-\theta\|^2\\
 &=d+\bigl(a^2-2a(d-2)\bigr)
                 \E_\theta\frac1{\|X\|^2}.
 \end{aligned}
\]
最后的期望有限且严格为正，所以对任意 $\theta$，最优常数同为 $a=d-2$。风险为
\[
 d-(d-2)^2\E_\theta\frac1{\|X\|^2}<d,
\]
而直接估计 $X$ 的风险为 $d$。这比较的是全部坐标的平方误差总和，并不声称每个坐标分别改善。
独立性及 $\E(Y-\theta)=0$ 使预测交叉项消失，所以 $\E\|Y-\delta_a(X)\|^2=d+R_\theta(a)$；最优收缩常数不变，最优预测风险严格小于直接预测 $X$ 所得的 $2d$。
\insight 允许适当偏差可以降低总平方风险。高斯分部积分将难处理的交叉误差改写为收缩函数的散度。
